\section{Comparaison entre l'étape 1 et l'étape 2}

\subsection{Comparaison des temps moyens}

\begin{figure}[H]
    \centering
    \begin{tikzpicture}
\begin{axis}[
    width=\linewidth,
    height=7.0cm,
    xlabel={Nombre de threads},
    ylabel={Temps moyen par pas (ms)},
    xmin=1,
    xmax=16,
    ymin=0,
    xtick={1,2,4,8,16},
    legend style={at={(0.5,-0.2)}, anchor=north, legend columns=1},
    grid=both,
    minor y tick num=1,
    axis line style={draw=black!60},
]
\addplot[
    color=enstaBleuFonce,
    mark=*,
    line width=1.1pt,
]
coordinates {
        (1, 227.970)
        (2, 124.934)
        (4, 69.480)
        (8, 48.834)
        (16, 33.233)
};
\addlegendentry{Étape 1 compute-only}

\addplot[
    color=enstaBleuClair,
    mark=square*,
    line width=1.1pt,
]
coordinates {
        (1, 246.802)
        (2, 138.732)
        (4, 79.649)
        (8, 52.546)
        (16, 42.668)
};
\addlegendentry{Étape 2 rank 1 compute}

\addplot[
    color=black!70,
    mark=triangle*,
    line width=1.1pt,
]
coordinates {
        (1, 246.919)
        (2, 138.853)
        (4, 79.767)
        (8, 52.673)
        (16, 42.801)
};
\addlegendentry{Étape 2 end-to-end}
\end{axis}
\end{tikzpicture}

    \caption{Comparaison des temps moyens par pas hors affichage}
    \label{fig:stage1_stage2_mean_times}
\end{figure}

Sur le benchmark sans fenêtre, l'étape 2 reste logiquement un peu plus coûteuse que l'étape 1, car elle ajoute la synchronisation MPI et l'envoi des positions. La courbe montre toutefois que l'écart entre \texttt{rank1\_compute} et \texttt{end-to-end} reste très faible : le surcoût MPI ne change pas l'ordre de grandeur du temps moyen par pas.

\subsection{Comparaison du speed-up compute-only}

\begin{figure}[H]
    \centering
    \begin{tikzpicture}
\begin{axis}[
    width=\linewidth,
    height=6.5cm,
    xlabel={Nombre de threads},
    ylabel={Speed-up},
    xmin=1,
    xmax=16,
    ymin=0,
    xtick={1,2,4,8,16},
    legend style={at={(0.5,-0.2)}, anchor=north, legend columns=2},
    grid=both,
    minor y tick num=1,
    axis line style={draw=black!60},
]
\addplot[
    color=enstaBleuFonce,
    mark=*,
    line width=1.1pt,
]
coordinates {
        (1, 1.000)
        (2, 1.825)
        (4, 3.281)
        (8, 4.668)
        (16, 6.860)
};
\addlegendentry{Étape 1 compute-only}

\addplot[
    color=black!70,
    mark=triangle*,
    line width=1.1pt,
]
coordinates {
        (1, 1.000)
        (2, 1.778)
        (4, 3.096)
        (8, 4.688)
        (16, 5.769)
};
\addlegendentry{Étape 2 end-to-end}
\end{axis}
\end{tikzpicture}

    \caption{Comparaison des speed-up compute-only entre l'étape 1 et l'étape 2}
    \label{fig:stage1_stage2_compute_speedup}
\end{figure}

L'étape 1 garde un léger avantage en speed-up compute-only. Par exemple, à 16 threads, elle atteint un speed-up de 6.861 contre 5.769 pour l'étape 2 en temps bout-en-bout. Cela reste cohérent : l'étape 2 ajoute une couche MPI, même si son coût de communication reste très faible.

\subsection{Comparaison interactive}

\begin{figure}[H]
    \centering
    \begin{tikzpicture}
\begin{axis}[
    ybar,
    bar width=14pt,
    width=\linewidth,
    height=6.8cm,
    ymin=0,
    ylabel={Temps moyen (ms)},
    symbolic x coords={Render, Update / attente},
    xtick=data,
    enlarge x limits=0.3,
    grid=both,
    minor y tick num=1,
    legend style={at={(0.5,-0.18)}, anchor=north, legend columns=2},
    axis line style={draw=black!60},
]
\addplot[fill=enstaBleuClair, draw=enstaBleuFonce] coordinates {
    (Render, 4.750)
    (Update / attente, 5.800)
};
\addlegendentry{Étape 1 interactive}

\addplot[fill=black!25, draw=black!70] coordinates {
    (Render, 4.425)
    (Update / attente, 6.025)
};
\addlegendentry{Étape 2 interactive}
\end{axis}
\end{tikzpicture}

    \caption{Décomposition des temps interactifs à 4 threads}
    \label{fig:stage1_stage2_interactive_breakdown}
\end{figure}

À 4 threads, les deux étapes ont pratiquement le même coût interactif total : \SI{10.550}{ms} pour l'étape 1 contre \SI{10.450}{ms} pour l'étape 2. L'étape 2 dépense un peu plus de temps en attente MPI, mais elle compense par un rendu légèrement plus court.

\begin{figure}[H]
    \centering
    \begin{tikzpicture}
\begin{axis}[
    width=\linewidth,
    height=6.5cm,
    xlabel={Nombre de threads},
    ylabel={Speed-up total},
    xmin=1,
    xmax=16,
    ymin=0,
    xtick={1,2,4,8,16},
    legend style={at={(0.5,-0.2)}, anchor=north, legend columns=2},
    grid=both,
    minor y tick num=1,
    axis line style={draw=black!60},
]
\addplot[
    color=enstaBleuFonce,
    mark=*,
    line width=1.1pt,
]
coordinates {
        (1, 1.000)
        (2, 1.503)
        (4, 2.019)
        (8, 2.427)
        (16, 2.506)
};
\addlegendentry{Étape 1 interactive}

\addplot[
    color=black!70,
    mark=triangle*,
    line width=1.1pt,
]
coordinates {
        (1, 1.000)
        (2, 1.789)
        (4, 2.256)
        (8, 2.449)
        (16, 2.840)
};
\addlegendentry{Étape 2 interactive}
\end{axis}
\end{tikzpicture}

    \caption{Comparaison des speed-up interactifs totaux entre l'étape 1 et l'étape 2}
    \label{fig:stage1_stage2_interactive_speedup}
\end{figure}

Le résultat le plus utile pour la suite est que l'étape 2 scale légèrement mieux en mode interactif. Le gain reste modéré, mais il existe : à 16 threads, on obtient un speed-up total de 2.840 contre 2.506 pour l'étape 1. Autrement dit, la séparation affichage/calcul n'apporte pas un saut massif immédiat, mais elle isole correctement le rendu sans pénaliser le coût par image, ce qui prépare proprement l'étape suivante.
